\documentclass[10pt,a4paper,sans]{moderncv}
\moderncvstyle{classic}
\moderncvcolor{blue}
\usepackage[utf8]{inputenc}
\usepackage[scale=0.91]{geometry}
\setlength{\hintscolumnwidth}{3cm}

% Personal data
\name{Dmitriy}{Kruglikov}
\address{Chemin des Moissons 8}{1294 Genthod, Suisse}
\phone[mobile]{+41783279082}
\email{dimitriy.kruglikov@gmail.com}
\social[linkedin]{dmitriy-kruglikov}
\social[github]{DimiK1337}
\photo[60pt][0.4pt]{../assets/images/DmitriyKruglikov_pfp.jpeg}

\begin{document}

\makecvtitle

\section{Profil}
Diplômé en informatique (Bachelor), de nationalités suisse, russe et américaine, avec une expérience en développement front-end et back-end, visualisation de données et encadrement pédagogique. Capacité démontrée à produire du code efficace et évolutif, avec une attention particulière à la performance et à la clarté. Intérêt marqué pour l’extraction d’insights à partir de données complexes.

\section{Expérience}
\cventry{Mar -- Juil 2024}{Ingénieur Front-End — Visualisation de données}{Effect Healthcare}{Amsterdam, Pays-Bas}{}{
\begin{itemize}
  \item Développement d’un \textit{clinical pathway mapper} avec \textbf{D3.js}, en optimisant les performances et l’utilisation mémoire.
  \item Recherche comparative : \textbf{POO vs programmation fonctionnelle}, \textbf{tuples vs objets}.
  \item Collaboration avec des développeurs seniors ; utilisation de \textbf{Wallaby} et \textbf{Karma} pour les tests.
\end{itemize}}

\cventry{Juil -- Août 2023}{Stagiaire Data Engineer}{Oxaigen}{Amsterdam, Pays-Bas}{}{
\begin{itemize}
  \item Développement d’une base de données PostgreSQL avec \textbf{Docker}, \textbf{PSQL} et \textbf{Bash}.
  \item Fusion et normalisation de schémas issus de plusieurs sources ; automatisation de l’initialisation de la base.
\end{itemize}}

\cventry{Mar -- Juin 2023}{Stagiaire Développeur Back-End}{Mozo}{Amsterdam, Pays-Bas}{}{
\begin{itemize}
  \item Développement back-end avec \textbf{Node.js}, \textbf{TypeScript} et \textbf{PostgreSQL}.
  \item Conception d’API REST et optimisation des opérations base de données lors d’un hackathon d’équipe.
\end{itemize}}

\cventry{2022 -- 2024}{Assistant d’enseignement — Programmes BSc Informatique \& IA}{Vrije Universiteit Amsterdam}{}{}{
\begin{itemize}
  \item Support sur plusieurs cours fondamentaux : \textbf{Introduction à la programmation (Java)}, \textbf{Advanced Programming}, \textbf{Requirements Engineering}, \textbf{Human-Computer Interaction [CS]}, \textbf{Human-Computer Interaction [AI]}.
  \item Accompagnement des étudiants sur les devoirs hebdomadaires, supervision de labs, retours techniques et design.
\end{itemize}}

\cventry{Juin 2019}{Stagiaire Data Analyst}{Office des Nations Unies à Genève}{}{Suisse}{
\begin{itemize}
  \item Automatisation de rapports mensuels avec \textbf{Python (pandas)} pour agréger et analyser des données CSV.
\end{itemize}}

\section{Formation}
\cventry{2021 -- 2024}{Bachelor (BSc) en informatique}{Vrije Universiteit Amsterdam}{}{}{}
\cventry{2019 -- 2021}{Diplôme du Baccalauréat International (IB)}{International School of Geneva}{}{}{}

\section{Projets}
\cvitem{Manga Downloader}{Application Python (GUI Kivy) pour télécharger et lire des mangas en anglais et en japonais (Android + PC).}
\cvitem{Kanji Study App}{Outil d’étude des kanjis (Japanese B, IB) développé avec Python/Kivy.}
\cvitem{\href{https://silveradocapitaladvisors.com/}{\textcolor{blue}{Site Silverado Capital Advisors}}}{Développement d’un site corporate en production avec \textbf{Next.js 14}, \textbf{TypeScript} et \textbf{Material UI}. Mise en place d’une architecture SEO scalable avec \textbf{Next Metadata API}, \textbf{Open Graph} et des données structurées \textbf{JSON-LD}.}

\section{Certificats}
\cvitem{2025}{\href{https://studies.cs.helsinki.fi/stats/api/certificate/fullstackopen/en/efb31cf5c06e1f21354b82294f6b1869}{Full Stack Open 2024 – Université d’Helsinki}}
\cvitem{2024}{\href{https://www.udemy.com/certificate/UC-4ab479b6-7be6-4212-880e-8d23d3724b80/}{Udemy 2024 – Learn and Understand D3.js for Data Visualization}}

\section{Compétences}
\cvitem{Langages}{Python, Java, TypeScript, JavaScript, SQL, C++, Bash}
\cvitem{Front-end}{React.js, Redux Toolkit, React Query, HTML/CSS, Vite}
\cvitem{Back-end}{Node.js, Express.js, API REST, JWT, PostgreSQL, MongoDB, Mongoose}
\cvitem{Outils}{Docker, Git, GitHub Workflows, Jest, Supertest, ESLint, dotenv, Wallaby, Karma}
\cvitem{Data}{pandas, matplotlib, BeautifulSoup, seaborn, scikit-learn}
\cvitem{Autres}{Unity (2D/3D), D3.js, Adobe Photoshop/Illustrator/XD, Kivy}

\end{document}